\subsection{自由同伦的圈}

定义 4. --- 设 X 是一个拓扑空间，c 和 c' 是 X 中的两个圈。若同伦 σ 连接 c 与 c'，且对所有 s ∈ I 都有 σ(0, s) = σ(1, s)，则称其为连接 c 与 c' 的自由同伦。如果存在连接 c 与 c' 的自由同伦，就称 c 与 c' 自由同伦。

连接 \(c\) 与 \(c'\) 的自由同伦对应于 X 的圈空间 \(\Omega(\mathrm{X})\) 中连接 \(c\) 与 \(c'\) 的道路。因此，关系“ \(c\) 与 \(c'\) 自由同伦”是 \(\Omega(\mathrm{X})\) 上的等价关系，其等价类是 \(\Omega(\mathrm{X})\) 的道路连通分支。

注. --- 设 \(\varphi\) 是从 \(\mathbf{R}\) 到 \(\mathbf{T} = \mathbf{R}/\mathbf{Z}\) 的典范映射（TG, V, p.~2）。映射 \(f \mapsto f \circ \varphi|\mathbf{I}\) 是从 \(\mathcal{C}_c(\mathbf{T};\mathbf{X})\) 到 \(\Omega(\mathbf{X})\) 的同胚，从而取道路连通分支后，得到从集合 {[}T; X{]}（III, p.~230）到 X 中的自由同伦圈类集合的双射。

命题 6. --- 设 X 是一个道路连通的拓扑空间，x 是 X 中的一点。

\begin{enumerate}
\def\labelenumi{\alph{enumi})}
\item
  X 中的每个圈都与 x 处的一个圈自由同伦。更确切地说，若 c 是 y 处的圈，d 是起点为 y、终点为 x 的道路，则 c 与 x 处的圈 \((\overline{d} * c) * d\) 自由同伦。
\item
  X 中 x 处的两个圈自由同伦，当且仅当它们的严格同伦类在群 \(\pi_{1}(X,x)\) 中共轭。
\end{enumerate}

我们来证明 \(a\)。对每个 \(s \in [0,1]\)，记 \(d_s\) 为 X 中由 \(d_s(t) = d(st)\)（\(t \in \mathbf{I}\)）定义的道路；其起点为 \(y\)。由于映射 \((s,t) \mapsto d(st)\) 连续，从 \(s \mapsto d_s\) 到 \(\mathbf{I}\) 的映射 \(\mathcal{C}_c(\mathbf{I};\mathbf{X})\) 是连续的（III, p.~257, 命题 1）。于是映射 \(s \mapsto (\overline{d_s} * c) * d_s\) 是 \(\Omega(\mathbf{X})\) 中的一条道路（III, p.~257, 命题 2），连接 \((e_y * c) * e_y\) 与 \((\overline{d} * c) * d\)，故得 \(a\)。

设 c 和 \(c'\) 是 X 中 x 处的两个圈。若它们的严格同伦类在 \(\pi_{1}(X,x)\) 中共轭，则存在 x 处的圈 d，使得 \(c'\) 与圈 \((\overline{d} * c) * d\) 严格同伦。由 a)，c 与 \(c'\) 自由同伦。反之，假设存在连接 c 与 \(\varphi\) 的自由同伦 \(c'\)。置 \(d(t) = \varphi(0,t)\)，则也有 \(d(t) = \varphi(1,t)\)，且 d 是 x 处的圈。由 III, p.~295 的引理 1，圈 \(c * d\) 与 \(d * c'\) 严格同伦。因此，c 与 \(c'\) 的严格同伦类在 \(\pi_{1}(X,x)\) 中共轭。

系理. --- 设 X 是一个道路连通的拓扑空间，x 是 X 中的一点。命题 6 使我们能够定义从 X 中的自由同伦圈类集合到 \(\pi_{1}(X,x)\) 中的共轭类集合的典范双射。

